\documentclass[hyperref={pdfpagelabels=false},table, handout]{beamer}

\input{lec_common}

\begin{document}

\part{Unit 1: First steps - A bit of everything.}

%\frame{\tableofcontents}
\begin{frame}{This Lecture}

  \begin{itemize}
    \item About this course
    \item About Python
    \item Data types and variables
    \item Import modules
    \item Lists
    \item For loops
    \item Elementary plotting
  \end{itemize}
\end{frame}

\begin{frame}{Why Python?}
Python is ...
\begin{itemize}
%\item{Easy to write and to read}
\item Free and open source
\item It is a \alert{scripting language} - an interpreted \structure{not a compiled} language
\item It is object oriented with modern exception handling, dynamic typing etc.
\item It has plenty of libraries among others the scientific ones:
\small linear algebra;  visualisation tools: plotting, image analysis; differential equations solving; symbolic computations; statistics
; etc.
\item It has possible usages: \small
Scientific computing, scripting, web sites, text parsing, data mining, ...
\end{itemize}
\end{frame}



\begin{frame}[fragile]{Premisses}
\begin{itemize}
\item Python version $\ge 3.x$ (installed through Anaconda).
\item \structure{IPython} shell
\item Work bench \structure{Spyder} (many others exist)
\item We work (later) with the Jupyter notebook.
\item All programs in this course start with the line
\begin{python}
from numpy import *
\end{python}
(will be explained later).
\end{itemize}
\end{frame}
\begin{frame}[fragile]{Course Book}
  Scientific Computing with Python 3 (2nd ed.)\\
Claus Führer, Jan Erik Solem, Olivier Verdier\\
Packt Publishing, 2021
\\[1cm]
\url{www.packtpub.com/}
\\[1cm]
\begin{center}
\includegraphics[width=2cm]{Python.jpg}
\end{center}
\end{frame}


\begin{frame}[fragile]{Examples}
Python may be used in \alert{interactive} mode \\
(executed in an \texttt{IPython} shell)
\begin{python}
In [2]: x = 3
In [3]: y = 5
In [4]: print(x + y)
8
\end{python}

Note:
\begin{python}
In [2]:
\end{python}
  is the \emph{prompt string} of the \texttt{IPython} shell. It counts the statements.
  In the more basic \texttt{Python} shell the \emph{prompt string} is \texttt{>\!>\!>}.
\end{frame}
\begin{frame}[fragile]{Examples: Linear Algebra}
Let us solve
\[\begin{pmatrix}1 &2\\3 &4\end{pmatrix} \cdot x = \begin{pmatrix}2\\1\end{pmatrix}\]
\begin{python}
In [5]: from scipy.linalg import solve
In [6]: M = array([[1., 2.],
                   [3., 4.]])
In [7]: V = array([2., 1.])
In [8]: x = solve(M, V)
In [9]: print(x)
[-3.   2.5]
\end{python}
Note: A line can be continued on the next line \emph{without any continuation symbol}
as long as parentheses or \emph{brackets are not closed}.
\end{frame}
\begin{frame}[fragile]{More examples}
Computing $\ee^{\ii\pi}$ and $2^{100}$:
\begin{python}
In [10]: print(exp(1j*pi)) # should return -1
(-1+1.22460635382e-16j)
In [11]: print(2**100)
1267650600228229401496703205376
\end{python}
Note: Everything following \pyth!#! is treated as a comment.\\
\vfill
{\tiny In this course expected results are displayed as comments directly after the command.}
\vfill
\end{frame}


\begin{frame}[fragile]{Numbers}
A number may be an \alert{integer}, a \alert{real number} or a \alert{complex number}.
The usual operations are
\begin{itemize}
\item \code{+} and \code{-} addition and substraction
\item \code{*} and \code{/} multiplication and division
\item \code{**} power
\end{itemize}
\begin{python}
2**(2+2) #  16
1.0j**2  # -1.0
\end{python}
\end{frame}

\begin{frame}[fragile]{Strings}
Strings are ``lists'' of characters, enclosed by simple or double quotes:
\begin{python}
str1 = 'valid string'
str2 = "string with double quotes"
\end{python}
\pause
You may also use \alert{triple quotes} for strings including multiple lines:
\begin{python}
str3 = """This is
a long,
long string"""
\end{python}
\end{frame}

\begin{frame}[fragile]{Import modules}
\begin{block}{Importing Numpy}
In order to use standard mathematical functions in Python, you must import some module.
 We will use \structure{numpy}.
\end{block}

\begin{python}
# You can import numpy and then
# use dot-notation after the name numpy
import numpy
print (numpy.exp(1j*numpy.pi))
# (-1+1.2246467991473532e-16j)
\end{python}

\pause

\begin{python}
# You can import numpy and rename the module
import numpy as np
print (np.exp(1j*np.pi))
# (-1+1.2246467991473532e-16j)
\end{python}

\pause

\begin{python}
# Or only import what you need
from numpy import exp, pi
print (exp(1j*pi))
# (-1+1.2246467991473532e-16j)
\end{python}

\end{frame}

\begin{frame}[fragile]{Variables}
\begin{block}{Variables}
A variable is a \alert{reference} to an object.
An object may have several references.
One uses the \alert{assignment operator} \code{=} to assign a value to a variable.
\end{block}

\vskip -5pt

\begin{example}
\begin{python}
i = 42      # integers
a = 3.14    # floats, the dot makes the difference
c = 5 + 1j  # complex numbers
sum = a + c # the sum will be a complex
print( a )         #  8.14 + 1j
print( type(sum) ) #  <class 'complex'>
\end{python}
\end{example}

\pause
\vskip -5pt

\begin{example}
\begin{python}
x = [3, 4] # a list object is created
y = x # this object is now referenced by x and by y
del x # we delete one of the references
del y # all references are deleted: the object is deleted
\end{python}
\end{example}
\end{frame}



\begin{frame}[fragile]{Concept: Lists}
\begin{block}{Lists}
A Python list is an \alert{ordered} list of objects, enclosed in square brackets. One accesses elements of a list using \alert{zero-based} indices inside square brackets.
\end{block}
\end{frame}

\begin{frame}[fragile]{List Examples}
\begin{example}
\begin{python}
L1 = [1, 2]
L1[0] # 1
L1[1] # 2
L1[2] # raises IndexError

L2 = ['a', 1, [3, 4]]
L2[0]    # 'a'
L2[2][0] # 3

L2[-1] # last element: [3,4]
L2[-2] # second to last: 1
\end{python}
\end{example}
\end{frame}

\begin{frame}[fragile]{List Utilities}
\begin{itemize}[<+->]
\item \pyth{range(n)} can be used to fill list with $n$ elements, starting with zero:
\begin{python}
list(range(5))  # [0, 1, 2, 3, 4]
\end{python}
\item \pyth{len(L)} gives the \alert{length} of a list:
\begin{python}
len(['a', 1, 2, 34]) # returns 4
\end{python}
\item Use \pyth{append} to append an element to a list:
\begin{python}
L = ['a', 'b', 'c']
L[-1] # 'c'
L.append('d')
L     # L is now ['a', 'b', 'c', 'd']
L[-1] # 'd'
\end{python}
\end{itemize}
\end{frame}

\begin{frame}[fragile]{List Comprehension}
A convenient way to build up lists is to use the \alert{list comprehension} construct, possibly with a conditional inside.
\begin{defblock}
The syntax of a list comprehension is
\begin{python}
[<expr> for <variable> in <list>]
\end{python}
\end{defblock}
\vfill
\begin{example}
\begin{python}
L  = [2, 3, 10, 1, 5]
L2 = [x*2 for x in L] # [4, 6, 20, 2, 10]
L3 = [x*2 for x in L if 4 < x <= 10] # [20, 10]
\end{python}
\end{example}
\end{frame}

\begin{frame}[fragile]{List Comprehension in Mathematics}
\begin{block}{Mathematical Notation}
This is very close to the mathematical notation for sets. Compare:
\[L_2 = \{2x;\ x\in L\}\]
and
\begin{python}
L2 = [2*x for x in L]
\end{python}
  One big difference though is, that lists are \emph{\alert{ordered}} while sets aren't.
\end{block}
\end{frame}

\begin{frame}[fragile]{Operations on Lists}
\begin{itemize}[<+->]
\item Adding two lists \alert{concatenates} (\emph{sammanfogar}) them:
\begin{python}
L1 = [1, 2]
L2 = [3, 4]
L = L1 + L2 # [1, 2, 3, 4]
\end{python}
\item Logically, multiplying a list with an integer concatenates the list with
itself several times: \texttt{n*L} is equivalent to $\underbrace{L+L+\cdots+L}_{\text{$n$ times}}$.
\vskip 8pt
\begin{python}
L = [1, 2]
3 * L # [1, 2, 1, 2, 1, 2]
\end{python}
(To multiply each element by $c$, we use \emph{arrays} instead of lists.)
\end{itemize}
\end{frame}



\begin{frame}[fragile]{\pyth!for! loops}
\begin{block}{for loop}
A \alert{for loop} allows to loop through a list using an \alert{index variable}. This variable takes succesively the values of the elements in the list.
\end{block}
\pause
\begin{example}
\begin{python}
L = [1, 2, 10]
for s in L:
  print(s * 2) # output: 2 4 20
\end{python}
\end{example}
\end{frame}

\begin{frame}[fragile]{Repeating a Task}
One \alert{typical use} of the \pyth{for} loop is to repeat a certain task a fixed number of times:
\begin{python}
n = 30
for i in range(n):
  do_something # this gets executed n times
\end{python}
\end{frame}


\begin{frame}[fragile]{Indentation}
The part to be repeated in the for loop has to be properly \alert{indented}:
\begin{python}
for elt in my_list:
  do_something()
  something_else()
  etc
print("loop finished") # outside the for block
\end{python}
\pause
\vfill
\alert{Note:} In contrast to other programming languages, the indentation in Python is \alert{mandatory}.
\vfill
Many other kinds of Python structures also have to be indented, we will cover this when introducing them.
\end{frame}


\begin{frame}[fragile]{Elementary Plotting}
We first make the central visualization tool in Python available:
\begin{python}
from matplotlib.pyplot import *
\end{python}

Then we generate two lists
\begin{python}
x_list=list(range(100)) # the first hundred numbers
y_list=[sqrt(x) for x in x_list]
\end{python}

And then we make a graph
\begin{python}
plot(x_list,y_list,'o')
title('My first plot')
xlabel('x')
ylabel('Square root of x')
show()
\end{python}
\end{frame}
\begin{frame}[fragile]{Elementary Plotting}
\includegraphics[width=.9\textwidth]{figure_1.png}
\end{frame}

\begin{frame}[fragile]{How to run a piece of code}
There are two phases in Python programming:
\begin{enumerate}
\item \alert{Writing} the code
\item \alert{Executing} (running) it
\end{enumerate}
\vfill
For Task 2, one needs to give the code to the Python \alert{interpreter} which reads the code and figures out what to do with it.
\vfill
Short snippets of code can be written directly in the interpreter and executed interactively. We use the interpreter \alert{IPython} which has extra features.
\vfill
For writing a larger program, one generally uses a text \alert{editor} which can highlight code in a good way.
\vfill
There are also development suites which bundle the interpreter, editor and other things into the same program.

\end{frame}

\begin{frame}[fragile]{Spyder}
Example of a development environment, Spyder
\includegraphics[height=.7\textheight]{Spyder}
\end{frame}

\begin{frame}[fragile]{IPython}
\vfill
IPython is an enhanced Python interpreter.
\vfill
You can launch it as a stand-alone application, but it is also embedded in Spyder.
\vfill
Some notes on usage:
\begin{itemize}
\item To execute the contents of a file named \code{myscript.py} just write \pyth{run myscript} in IPython
or use the green "run" arrow  \hspace{1em} \includegraphics[scale=0.6]{Spyder_run} \hspace{1em} in the Spyder toolbar.
\item Use the arrow keys\hspace{1em}\includegraphics[scale=0.2]{up_down_arrow_keys.jpg}\hspace{1em} to visit previous commands and
\item Use the tabulation key
\hspace{1em}\includegraphics[scale=0.2]{tab_key.jpg}\hspace{1em}
to auto-complete commands and names
\item To get \alert{help} on an object just type \code{?} after it and then return
\item When you want to quit, write \code{exit()}
\end{itemize}

\end{frame}



\end{document}
